\documentclass[twocolumn]{notes}
\usepackage{notes}

\begin{document}
\title{ChatGPT: solving integral}
\date{\today} % \formatdate{10}{11}{2022}
\maketitle
\tableofcontents
\newpage

\section{Methods for integrals of the form $\int_{\mathbb R^2} e^{-\beta V(x)}\,G(x-x_0)\,dx$}

Consider
\[
I(x_0)=\int_{\mathbb R^2} \exp\!\big(-\beta V(x)\big)\,G(x-x_0)\,dx,
\]
where $x\in\mathbb R^2$ and
\[
V(x)=\sum_{n=1}^N w_n\,G_n(x-x_n)
\]
is a sum of Gaussians (possibly with different amplitudes and widths). In general there is no closed form because this is an \emph{exponential of a sum of Gaussians}.

Below are practical methods beyond brute-force quadrature on an $x$--grid.

\subsection{View it as an expectation under $G$ (Monte Carlo / Quasi--Monte Carlo)}
If $G(\cdot-x_0)$ is a \emph{normalized} Gaussian density, e.g.
\[
G(x-x_0)=\mathcal N(x;\,x_0,\Sigma_0),
\]
then
\[
I(x_0)=\mathbb E_{X\sim \mathcal N(x_0,\Sigma_0)}\!\left[\exp\!\big(-\beta V(X)\big)\right].
\]
A simple estimator is
\[
I(x_0)\approx \frac{1}{K}\sum_{k=1}^K \exp\!\big(-\beta V(X_k)\big),
\qquad X_k\sim \mathcal N(x_0,\Sigma_0).
\]

Variance reduction (often very effective in 2D):
\begin{itemize}
\item Quasi--Monte Carlo (Sobol/Halton) in 2D (map low-discrepancy points through the Gaussian inverse CDF).
\item Antithetic sampling: sample $Z$ and $-Z$ in the latent standard normal, then set $X=x_0+LZ$.
\item Control variates using analytic approximations (see cumulant expansion below).
\end{itemize}

\subsection{Laplace (saddle-point) approximation}
Write the integrand as $\exp(-\Phi(x))$ where
\[
\Phi(x)=\beta V(x)-\log G(x-x_0).
\]
Then
\[
I(x_0)=\int_{\mathbb R^2} e^{-\Phi(x)}\,dx.
\]
If $\Phi$ has a dominant minimizer $x^*$ (mode of the integrand), the Laplace approximation gives
\[
I(x_0)\approx (2\pi)\,\frac{\exp\!\big(-\Phi(x^*)\big)}{\sqrt{\det H}},
\qquad H=\nabla^2\Phi(x^*).
\]
Notes:
\begin{itemize}
\item In 2D, finding $x^*$ by Newton or quasi-Newton is typically cheap.
\item $\nabla V$ and $\nabla^2 V$ are analytic for Gaussians, so $H$ is accessible.
\item If there are multiple comparable minima, one can use a multi-modal Laplace approximation by summing contributions from several stationary points.
\end{itemize}

\subsection{Cumulant / perturbation expansion under the Gaussian $G$}
Again let $X\sim G(\cdot-x_0)$ (normalized). Then
\[
I(x_0)=\mathbb E\big[e^{-\beta V(X)}\big].
\]
If $V(X)$ does not fluctuate too wildly under $G$, expand $\log I$ in cumulants:
\[
\log I \approx -\beta \mu_V + \frac{\beta^2}{2}\sigma_V^2 - \frac{\beta^3}{6}\kappa_3+\cdots,
\]
where under the distribution of $X$,
\[
\mu_V=\mathbb E[V(X)],\qquad \sigma_V^2=\mathrm{Var}(V(X)),\qquad \kappa_3=\text{3rd cumulant of }V(X).
\]
Key point for your structure: since $V$ is a sum of Gaussians, quantities like
\[
\mathbb E\big[G_n(X-x_n)\big]
\quad\text{and}\quad
\mathbb E\big[G_n(X-x_n)\,G_m(X-x_m)\big]
\]
are analytic (Gaussian integrals involving products/convolutions of Gaussians). Therefore $\mu_V$ and $\sigma_V^2$ can often be computed in closed form, enabling a purely analytic 2nd-order approximation:
\[
I(x_0)\approx \exp\!\left(-\beta \mu_V + \frac{\beta^2}{2}\sigma_V^2\right).
\]
This approximation can also be used as a control variate to reduce Monte Carlo variance.

\subsection{Importance sampling with an adaptive Gaussian proposal}
If sampling from $G$ is inefficient (e.g.\ large $\beta$ makes the integrand sharply concentrated away from $x_0$), approximate the target
\[
\pi(x)\propto e^{-\beta V(x)}G(x-x_0)
\]
by a Gaussian proposal $q(x)=\mathcal N(m,S)$, for instance from a Laplace fit around the mode ($m=x^*$, $S=H^{-1}$). Then
\[
I(x_0)=\int e^{-\beta V(x)}G(x-x_0)\,dx
= \mathbb E_{X\sim q}\left[\frac{e^{-\beta V(X)}G(X-x_0)}{q(X)}\right],
\]
which can reduce variance substantially.

\subsection{Convolution/Fourier viewpoint (if you need $I(x_0)$ for many $x_0$)}
Define
\[
f(x)=e^{-\beta V(x)}.
\]
Then
\[
I(x_0)=\int f(x)\,G(x-x_0)\,dx = (f*G)(x_0),
\]
i.e.\ $I$ is a Gaussian smoothing of $f$. If you need $I(x_0)$ on a grid of $x_0$ values, you can:
\begin{itemize}
\item compute $f$ on a grid once,
\item convolve with $G$ using FFT: $\widehat{I}(k)=\widehat{f}(k)\widehat{G}(k)$.
\end{itemize}
This is still grid-based, but it avoids performing a separate 2D integration for each $x_0$.

\subsection{Gauss--Hermite quadrature after a Gaussian change of variables}
If $G(x-x_0)=\mathcal N(x_0,\Sigma_0)$, set $x=x_0+Lz$ where $LL^\top=\Sigma_0$ and $z\sim\mathcal N(0,I)$. Then
\[
I(x_0)=\int_{\mathbb R^2} \exp\!\big(-\beta V(x_0+Lz)\big)\,\varphi(z)\,dz,
\]
where $\varphi$ is the standard normal density in $\mathbb R^2$. This is a Gaussian-weighted integral, where:
\begin{itemize}
\item tensor-product Gauss--Hermite quadrature can converge quickly in 2D for smooth integrands,
\item sparse-grid (Smolyak) variants can reduce the number of nodes.
\end{itemize}

\subsection{What to choose in practice}
\begin{itemize}
\item Few $x_0$ evaluations, moderate/large $\beta$: Laplace approximation; optionally add an importance-sampled Monte Carlo correction.
\item Moderate $\beta$ and $V(X)$ not too variable under $G$: 2nd-order cumulant approximation (analytic mean/variance), possibly combined with Monte Carlo control variates.
\item Many $x_0$ values on a grid: compute $f=e^{-\beta V}$ on the grid once and convolve with $G$ via FFT.
\item Robust black-box in 2D without grids: Quasi--Monte Carlo sampling under $G$.
\end{itemize}

\subsection{Two clarifying questions (to pick the best method)}
\begin{enumerate}
\item Is $G(x-x_0)$ a normalized Gaussian density (with covariance $\Sigma_0$)? Are the Gaussians inside $V$ normalized or unnormalized ``bumps''?
\item Are the weights $w_n$ mostly positive, or can they be negative? What is a typical scale for $\beta w_n$?
\end{enumerate}


\end{document}



%% \href{https://www.ferglab.com/}{\color{darkred}Website}

% \usepackage{minted}
% \begin{minted}[mathescape, linenos]{python}
%     # Note: $\pi=\lim_{n\to\infty}\frac{P_n}{d}$
%     title = "Hello World"
% \end{minted}

%%=======================================================================
%%     PREFERRED FIGURE/TABLE
%%=======================================================================
%%                                 %%-------------------------------------%%
%% \begin{figure}[h!]              %%  caption right of figure
%%   \begin{minipage}[c]{0.67\textwidth}
%%     \includegraphics[width=\textwidth]{imgfile.png}
%%   \end{minipage}\hfill
%%   \begin{minipage}[c]{0.3\textwidth}
%%     \caption{Caption on the right hand side}
%%     \label{fig:}
%%   \end{minipage}
%% \end{figure}
%%                                 %%-------------------------------------%%
%% \begin{figure}[h!]\centering    %%  caption below figure
%%   \includegraphics[width=0.5\textwidth]{imgfile.png}
%%   \caption{General caption}
%%   \label{fig:label}
%% \end{figure}
%%                                 %%-------------------------------------%%
%% \begin{figure}[!h]\centering    %%  array of figures with subcaptions
%%   \subfloat[]{\includegraphics[width=0.5\textwidth]{}}
%%   \subfloat[]{\includegraphics[width=0.5\textwidth]{}} \\
%%   \subfloat[]{\includegraphics[width=0.5\textwidth]{}}
%%   \subfloat[]{\includegraphics[width=0.5\textwidth]{}} \\
%%   \subfloat[]{\includegraphics[width=0.5\textwidth]{}}
%%   \subfloat[]{\includegraphics[width=0.5\textwidth]{}}
%%   \caption{General caption}
%%   \label{fig:label}
%% \end{figure}
%%                                        %%------------------------------%%
%% \begin{table}[!h]\centering\ra{1.3}    %%  simple table
%%   \begin{tabular}{@{}rrrr@{}}\toprule
%%     & $t=0$ & $t=1$ & $t=2$ \\ \midrule
%%     $c$ & 1 & 2 & 3 \\
%%     $c$ & 1 & 2 & 3 \\
%%     \bottomrule
%%   \end{tabular}
%%   \caption{Caption}
%% \end{table}
%%                                        %%------------------------------%%
%% \begin{table}[!h]\centering\ra{1.3}    %%  complicated table
%%   \begin{tabular}{@{}rrrrcrrr@{}}\toprule
%%     & \multicolumn{3}{c}{$w = 8$} & \phantom{abc}&
%%     \multicolumn{3}{c}{$w = 16$}\\
%%     \cmidrule{2-4} \cmidrule{6-8} \cmidrule{10-12}
%%     & $t=0$ & $t=1$ & $t=2$ && $t=0$ & $t=1$ & $t=2$ \\ \midrule
%%     \rowcolor[gray]{.95} $c$ & 1 & 2 & 3 && 4 & 5 & 6 \\
%%     \rowcolor[gray]{1.0} $c$ & 1 & 2 & 3 && 4 & 5 & 6 \\
%%     \bottomrule
%%   \end{tabular}
%%   \caption{Caption}
%% \end{table}

%%=======================================================================
%%     OLDER FIGURE/TABLE
%%=======================================================================
%% \begin{center}
%%   \includegraphics[width=0.75\linewidth]{}
%%   \captionof{figure}{}
%% \end{center}
%%
%% \begin{center}
%%   \begin{tabular}{rl}
%%     \includegraphics[width=0.5\linewidth]{} &
%%     \includegraphics[width=0.5\linewidth]{} \\
%%     \includegraphics[width=0.5\linewidth]{} &
%%     \includegraphics[width=0.5\linewidth]{} \\
%%   \end{tabular}
%%   \captionof{figure}{}
%% \end{center}
%%
%% \begin{center}
%%   \begin{tabular}{lll}
%%     \hline
%%     A & B & C \\
%%     \hline
%%     \hline
%%     1 & 2 & 3 \\
%%     \hline
%%   \end{tabular}
%%   \captionof{table}{test}
%% \end{center}
